\documentclass[11pt,a4paper]{article}
% ============================================================
% Préambule commun -- Module M114
% Mathématiques appliquées à la gestion
% Licence MGE -- Université Mundiapolis Business School
% ============================================================
\usepackage[a4paper,top=2.9cm,bottom=2.5cm,left=2.3cm,right=2.3cm,headheight=1.85cm,headsep=8pt]{geometry}
\usepackage[utf8]{inputenc}
\usepackage[T1]{fontenc}
\usepackage{lmodern}
\usepackage[french]{babel}
\usepackage{amsmath,amssymb,mathtools,amsthm}
\usepackage{graphicx}
\usepackage[export]{adjustbox}
\usepackage[dvipsnames]{xcolor}
\usepackage{titlesec}
\usepackage{enumitem}
\usepackage{fancyhdr}
\usepackage{tcolorbox}
\tcbuselibrary{skins,breakable}
\usepackage{array}
\usepackage{booktabs}
\usepackage{tabularx}
\usepackage{multirow}
\usepackage{tikz}
\usetikzlibrary{arrows.meta,calc,positioning}
\usepackage{csquotes}
\usepackage[hidelinks]{hyperref}

% ---------- Couleurs Mundiapolis ----------
\definecolor{mundiblue}{HTML}{0070C0}
\definecolor{mundidark}{HTML}{123B5E}
\definecolor{mundigray}{HTML}{5A5A5A}
\definecolor{munditeal}{HTML}{1B6B6B}
\definecolor{mundired}{HTML}{B03A2E}

% ---------- Titres de sections ----------
\titleformat{\section}
  {\normalfont\Large\bfseries\color{mundidark}}
  {\thesection}{0.6em}{}
\titleformat{\subsection}
  {\normalfont\large\bfseries\color{mundiblue}}
  {\thesubsection}{0.6em}{}
\titleformat{\subsubsection}
  {\normalfont\normalsize\bfseries\color{mundigray}}
  {\thesubsubsection}{0.6em}{}

% ---------- Logos Honoris (gauche) / Mundiapolis (droite) ----------
\graphicspath{{logos/}{./}}
\newcommand{\logoheight}{1.35cm}
% Même hauteur + valign=c : les deux logos sont calés sur le même axe horizontal.
% On ne fixe PAS la largeur : les formats d'origine (Honoris / Mundiapolis) sont différents.
\newcommand{\putlogo}[1]{%
  \includegraphics[height=\logoheight,keepaspectratio,valign=b]{#1}%
}
\newcommand{\headerlogos}{%
  \noindent\makebox[\textwidth][s]{%
    \putlogo{honoris.jpg}\hfill
    \putlogo{munidia.png}%
  }%
}

% ---------- En-tête / pied de page ----------
\pagestyle{fancy}
\fancyhf{}
\renewcommand{\headrulewidth}{0pt}
\renewcommand{\footrulewidth}{0.4pt}
\fancyhead[L]{\putlogo{honoris.jpg}}
\fancyhead[R]{\putlogo{munidia.png}}
\fancyfoot[L]{\small\color{mundigray} Mathématiques appliquées à la gestion -- S1 MGE}
\fancyfoot[C]{\small\color{mundigray}\thepage}
\fancyfoot[R]{\small\color{mundigray} Pr.\ Rajae Malek}
\fancypagestyle{plain}{%
  \fancyhf{}%
  \renewcommand{\headrulewidth}{0pt}%
  \renewcommand{\footrulewidth}{0.4pt}%
  \fancyhead[L]{\putlogo{honoris.jpg}}%
  \fancyhead[R]{\putlogo{munidia.png}}%
  \fancyfoot[C]{\small\color{mundigray}\thepage}%
}

% ---------- Boîtes pédagogiques ----------
\newtcolorbox{definitionbox}[1][]{
  breakable, enhanced, colback=blue!4, colframe=mundiblue,
  boxrule=0.6pt, arc=1mm, left=6pt, right=6pt, top=4pt, bottom=4pt,
  fonttitle=\bfseries, coltitle=white, colbacktitle=mundiblue,
  title={Définition #1}
}
\newtcolorbox{propbox}[1][]{
  breakable, enhanced, colback=green!5, colframe=OliveGreen!70!black,
  boxrule=0.6pt, arc=1mm, left=6pt, right=6pt, top=4pt, bottom=4pt,
  fonttitle=\bfseries, coltitle=white, colbacktitle=OliveGreen!70!black,
  title={Propriété #1}
}
\newtcolorbox{theoremebox}[1][]{
  breakable, enhanced, colback=blue!6, colframe=mundidark,
  boxrule=0.8pt, arc=1mm, left=6pt, right=6pt, top=4pt, bottom=4pt,
  fonttitle=\bfseries, coltitle=white, colbacktitle=mundidark,
  title={Théorème #1}
}
\newtcolorbox{exemplebox}[1][]{
  breakable, enhanced, colback=yellow!8, colframe=orange!80!black,
  boxrule=0.6pt, arc=1mm, left=6pt, right=6pt, top=4pt, bottom=4pt,
  fonttitle=\bfseries, coltitle=black, colbacktitle=orange!25,
  title={Exemple #1}
}
\newtcolorbox{gestionbox}[1][Application en gestion]{
  breakable, enhanced, colback=orange!6, colframe=orange!80!black,
  boxrule=0.6pt, arc=1mm, left=6pt, right=6pt, top=4pt, bottom=4pt,
  fonttitle=\bfseries, coltitle=black, colbacktitle=orange!30,
  title=#1
}
\newtcolorbox{exobox}[1][]{
  breakable, enhanced, colback=gray!5, colframe=mundidark,
  boxrule=0.6pt, arc=1mm, left=6pt, right=6pt, top=4pt, bottom=4pt,
  fonttitle=\bfseries, coltitle=white, colbacktitle=mundidark,
  title={Exercice #1}
}
\newtcolorbox{corrigebox}[1][]{
  breakable, enhanced, colback=gray!3, colframe=gray!60!black,
  boxrule=0.6pt, arc=1mm, left=6pt, right=6pt, top=4pt, bottom=4pt,
  fonttitle=\bfseries, coltitle=white, colbacktitle=gray!60!black,
  title={Corrigé -- #1}
}
\newtcolorbox{methodebox}[1][]{
  breakable, enhanced, colback=teal!4, colframe=munditeal,
  boxrule=0.6pt, arc=1mm, left=6pt, right=6pt, top=4pt, bottom=4pt,
  fonttitle=\bfseries, coltitle=white, colbacktitle=munditeal,
  title={Méthode #1}
}
\newtcolorbox{attentionbox}[1][Piège classique]{
  breakable, enhanced, colback=red!4, colframe=mundired,
  boxrule=0.6pt, arc=1mm, left=6pt, right=6pt, top=4pt, bottom=4pt,
  fonttitle=\bfseries, coltitle=white, colbacktitle=mundired,
  title=#1
}
\newtcolorbox{planbox}[1][Plan de la séance (3 heures)]{
  breakable, enhanced, colback=mundidark!4, colframe=mundidark,
  boxrule=0.6pt, arc=1mm, left=6pt, right=6pt, top=4pt, bottom=4pt,
  fonttitle=\bfseries, coltitle=white, colbacktitle=mundidark,
  title=#1
}
\newtcolorbox{objectifbox}[1][Objectifs]{
  breakable, enhanced, colback=mundiblue!4, colframe=mundiblue,
  boxrule=0.6pt, arc=1mm, left=6pt, right=6pt, top=4pt, bottom=4pt,
  fonttitle=\bfseries, coltitle=white, colbacktitle=mundiblue,
  title=#1
}
\newtcolorbox{remarquebox}[1][Remarque]{
  breakable, enhanced, colback=gray!6, colframe=mundigray,
  boxrule=0.5pt, arc=1mm, left=6pt, right=6pt, top=4pt, bottom=4pt,
  fonttitle=\bfseries, coltitle=white, colbacktitle=mundigray,
  title=#1
}

% ---------- Macros ----------
\newcommand{\N}{\mathbb{N}}
\newcommand{\R}{\mathbb{R}}
\newcommand{\un}{(u_n)}
\newcommand{\vn}{(v_n)}
\newcommand{\wn}{(w_n)}
\newcommand{\limn}{\lim_{n\to +\infty}}
\newcommand{\MAD}{\,\mathrm{MAD}}
\newcommand{\eps}{\varepsilon}

\DeclareMathOperator{\card}{card}

\setlength{\parindent}{0pt}
\setlength{\parskip}{4pt}
\renewcommand{\arraystretch}{1.25}
\setlist[enumerate]{leftmargin=1.4cm, itemsep=3pt}
\setlist[itemize]{leftmargin=1.2cm, itemsep=2pt}

% Page de garde générique
% \pagedegarde{sous-titre}{nature du document}
\newcommand{\pagedegarde}[2]{%
\begin{titlepage}
\hypersetup{pageanchor=false}
\headerlogos\\[0.7cm]
\centering
{\Huge\bfseries\color{mundiblue} UNIVERSITÉ MUNDIAPOLIS}\\[0.45cm]
{\color{mundigray}\large Business School}\\[0.25cm]
{\color{mundigray}\rule{0.5\textwidth}{0.4pt}}\\[1.1cm]
{\Large\color{mundidark} #2}\\[0.45cm]
{\Huge\bfseries\color{mundiblue} Mathématiques appliquées\\[0.15cm] à la gestion}\\[0.7cm]
{\Large\itshape\color{mundidark} #1}\\[1.4cm]
\begin{tabular}{@{}r l@{}}
\textbf{Module :} & M114 \\[4pt]
\textbf{Filière :} & \begin{tabular}[t]{@{}l@{}}Licence en Management\\ et Gestion des Entreprises (MGE)\end{tabular} \\[10pt]
\textbf{Niveau -- Semestre :} & S1 -- 1\up{ère} année \\[4pt]
\textbf{Enseignant(e) :} & Pr.\ Rajae Malek \\[4pt]
\textbf{Année universitaire :} & 2026 -- 2027 \\
\end{tabular}
\vfill
{\color{mundigray}\rule{0.7\textwidth}{0.4pt}}\\[0.3cm]
{\small\color{mundigray} Université Mundiapolis -- Tous droits réservés}
\end{titlepage}
\hypersetup{pageanchor=true}
}


\begin{document}

\pagedegarde{Chapitre 1 -- Les suites réelles}{Travaux dirigés -- Séance de 3 heures}

\section*{Consignes}

\begin{itemize}
  \item Durée : \textbf{3 heures}. Calculatrice scientifique autorisée. Justifier chaque réponse.
  \item Les exercices sont conçus pour être traités \textbf{dans l'ordre}. L'exercice 6 est un complément si le temps le permet.
  \item Les applications de gestion (exercice 5) constituent le cœur de la séance : y consacrer le temps indiqué.
\end{itemize}

\begin{planbox}[Déroulement indicatif]
\renewcommand{\arraystretch}{1.15}
\begin{tabularx}{\textwidth}{@{}c c X@{}}
\toprule
\textbf{Durée} & \textbf{Exercice} & \textbf{Compétence visée} \\
\midrule
25 min & 1 & Monotonie, bornes, théorème de la limite monotone. \\
20 min & 2 & Calculs de limites, formes indéterminées. \\
25 min & 3 & Gendarmes, suites extraites, critère de divergence. \\
30 min & 4 & Suite récurrente d'ordre 1 (démarche complète). \\
45 min & 5 & Arithmétique, géométrique, arithmético-géométrique en gestion. \\
25 min & 6 & Suites adjacentes \emph{ou} récurrence d'ordre 2. \\
10 min & --- & Mise en commun / synthèse. \\
\bottomrule
\end{tabularx}
\end{planbox}

% ------------------------------------------------------------
\section*{Exercice 1 -- Monotonie, bornes et convergence \hfill {\normalsize\normalfont (25 min -- 4 points)}}

On considère la suite $\un$ définie par
\[
u_n = \dfrac{3n+1}{n+2}.
\]

\begin{enumerate}
  \item Calculer $u_0$, $u_1$ et $u_2$.
  \item Calculer $u_{n+1}-u_n$ et en déduire le sens de variation de $\un$.
  \item Montrer que, pour tout $n\in\N$, $u_n < 3$. La suite est-elle minorée ?
  \item En déduire que $\un$ converge, puis calculer sa limite de deux façons :
  \begin{enumerate}
    \item en utilisant le théorème des opérations sur les limites ;
    \item en utilisant le théorème de la limite monotone (on identifiera la limite comme un majorant \og le plus petit possible \fg{}).
  \end{enumerate}
\end{enumerate}

% ------------------------------------------------------------
\section*{Exercice 2 -- Calculs de limites \hfill {\normalsize\normalfont (20 min -- 3 points)}}

Déterminer les limites suivantes (en justifiant, notamment en cas de forme indéterminée) :

\begin{enumerate}
  \item $\displaystyle\limn \dfrac{4n^3 - n + 2}{2n^3 + 5n - 1}$
  \item $\displaystyle\limn \bigl(\sqrt{n^2 + n} - n\bigr)$
  \item $\displaystyle\limn \dfrac{2^n + 3^n}{3^n + 1}$
\end{enumerate}

\begin{methodebox}[Rappels utiles]
\begin{itemize}
  \item Pour un quotient de polynômes : factoriser par le terme de plus haut degré.
  \item Pour une expression du type $\sqrt{a_n}-b_n$ : multiplier par la \textbf{quantité conjuguée}.
  \item Pour un quotient d'exponentielles : factoriser par le terme \textbf{dominant} (ici $3^n$).
\end{itemize}
\end{methodebox}

% ------------------------------------------------------------
\section*{Exercice 3 -- Gendarmes et suites extraites \hfill {\normalsize\normalfont (25 min -- 3,5 points)}}

\begin{enumerate}
  \item Soit $u_n = \dfrac{(-1)^n}{n+1} + \dfrac{2n}{n+3}$.
  \begin{enumerate}
    \item Encadrer $\dfrac{(-1)^n}{n+1}$ par deux suites convergeant vers $0$.
    \item En déduire $\limn u_n$.
  \end{enumerate}

  \item On pose $v_n = (-1)^n \left(1+\dfrac{1}{n}\right)$ pour $n\geq 1$.
  \begin{enumerate}
    \item Étudier les suites extraites $(v_{2n})$ et $(v_{2n+1})$.
    \item Conclure quant à la convergence de $\vn$.
  \end{enumerate}

  \item Soit $w_n = \dfrac{\sin n}{n+1}$. Montrer que $w_n \to 0$ (on rappelle que $\lvert\sin n\rvert \leq 1$).
\end{enumerate}

% ------------------------------------------------------------
\section*{Exercice 4 -- Suite récurrente \hfill {\normalsize\normalfont (30 min -- 4 points)}}

On considère la suite définie par $u_0 = 1$ et
\[
u_{n+1} = \sqrt{3u_n + 4}.
\]
On note $f$ la fonction $f(x)=\sqrt{3x+4}$, définie sur $\bigl[-\tfrac{4}{3},+\infty\bigr[$.

\begin{enumerate}
  \item Résoudre $f(\ell)=\ell$. Quelles sont les valeurs \emph{candidates} à être limite de $\un$ ?
  \item Montrer par récurrence que, pour tout $n\in\N$, $1 \leq u_n \leq 4$.
  \item Étudier le signe de $f(x)-x$ sur $[1,4]$ et en déduire que $\un$ est croissante.
  \item Conclure : $\un$ converge, et préciser sa limite.
\end{enumerate}

% ------------------------------------------------------------
\section*{Exercice 5 -- Suites particulières et applications en gestion \hfill {\normalsize\normalfont (45 min -- 5 points)}}

Cet exercice est le cœur de la séance. Les trois parties sont indépendantes.

\subsection*{A. Amortissement linéaire (suite arithmétique)}

Une machine-outil est achetée $80\,000\MAD$ et amortie linéairement de $8\,000\MAD$ par an. On note $V_n$ sa valeur comptable après $n$ années.

\begin{enumerate}
  \item Justifier que $(V_n)$ est une suite arithmétique. Préciser $V_0$ et la raison.
  \item Exprimer $V_n$ en fonction de $n$.
  \item Au bout de combien d'années la machine est-elle totalement amortie ?
  \item Calculer la valeur comptable cumulée des amortissements après $5$ ans.
\end{enumerate}

\subsection*{B. Capitalisation (suite géométrique)}

Un capital $C_0 = 15\,000\MAD$ est placé à intérêts composés au taux annuel $t=4{,}5\%$. On note $C_n$ le capital acquis après $n$ années.

\begin{enumerate}
  \item Exprimer $C_n$ en fonction de $n$.
  \item Calculer $C_8$ (on donnera la valeur exacte, puis une valeur arrondie à l'unité).
  \item Au bout de combien d'années le capital aura-t-il \textbf{doublé} ? (on pourra utiliser $\ln$).
\end{enumerate}

\subsection*{C. Emprunt à annuités constantes (suite arithmético-géométrique)}

Une entreprise emprunte $D_0 = 200\,000\MAD$ au taux annuel $t=4\%$, remboursé par une annuité constante $A=24\,000\MAD$. On note $D_n$ le capital restant dû \textbf{après} le $n$-ième versement.

\begin{enumerate}
  \item Justifier la relation $D_{n+1} = 1{,}04\, D_n - 24\,000$.
  \item Déterminer le point fixe $\ell$ de cette récurrence, puis exprimer $D_n$ en fonction de $n$.
  \item Calculer $D_5$ (arrondi à l'unité).
  \item Déterminer le plus petit entier $n$ tel que $D_n \leq 0$. Interpréter.
\end{enumerate}

% ------------------------------------------------------------
\section*{Exercice 6 -- Complément \hfill {\normalsize\normalfont (25 min -- 0,5 point de bonus)}}

Traiter \textbf{l'une} des deux questions, au choix.

\subsection*{Option A -- Suites adjacentes}

On pose, pour $n\geq 0$,
\[
u_n = \dfrac{n}{n+1}, \qquad v_n = \dfrac{n+2}{n+1}.
\]
Montrer que $\un$ et $\vn$ sont adjacentes. En déduire qu'elles convergent vers une même limite $\ell$, que l'on déterminera. Encadrer $\ell$ à $10^{-2}$ près à l'aide d'un rang bien choisi.

\subsection*{Option B -- Demande à mémoire de deux périodes}

La demande $D_n$ d'un produit vérifie
\[
D_{n+2} = 0{,}6\, D_{n+1} + 0{,}4\, D_n,
\]
avec $D_0 = 100$ et $D_1 = 120$ (en milliers d'unités).
\begin{enumerate}
  \item Écrire l'équation caractéristique et déterminer ses racines.
  \item En déduire l'expression de $D_n$ en fonction de $n$.
  \item Déterminer $\limn D_n$ et interpréter économiquement.
\end{enumerate}

\vspace{0.6cm}
\begin{center}
{\small\color{mundigray} Fin du TD -- Les corrigés détaillés sont dans \texttt{Chapitre1\_Suites\_reelles\_Corrections.pdf}}
\end{center}

\end{document}
